\subsection{Autonomous coherent-memory hierarchies}
Zonnios--Binder \cite[Proposition~2]{ZB90} prove monotonicity of their
memory-parametrized process distinguishability and its upper bound by
the strategy norm.  Their Theorem~1 compiles any finite-time tester into
repeated use of one instrument with an internal counter and sufficient
coherent memory; Corollary~1 gives finite-time completeness.  Their
unknown object is a multi-time process, and their resource is the coherent
ancilla dimension of an autonomous repeated-instrument machine.
Our reset class instead allows public time-dependent history rules and
unrestricted receiver storage, but forbids coherent import across a
specified fresh-preparation cut.  Its hard square cost $Q$ is not an
ancilla dimension.  Neither their completeness theorem nor our fixed-tube
profile law supplies the other's conclusion.  The exact finite-ensemble
hierarchy of Theorem~\ref{thm:operationalhierarchy90} likewise compares
three acquisition/receiver classes, not every dimension slice of their
hierarchy.  Their work is a direct antecedent for systematic coherent-memory
interpolation; no first such framework is claimed here.

\subsection{Classical outputs and accessible post-measurement states}
Eid--Quintino \cite{EQ90} consider the quantum state returned by the
associated L\"uders instrument in addition to the classical label.
Our unknown channel returns only that label.  The references used in
Theorem~\ref{thm:operationalhierarchy90} were supplied by the experimenter;
they are not residual quantum outputs of the device.  Thus an instrument
discrimination advantage cannot be imported into the present interface.

\subsection{What the exact hierarchy adds}
The strict tester inclusion in Proposition~\ref{prop:memorycomparison89}
alone is not an operational score theorem.  Theorem~\ref{thm:operationalhierarchy90}
provides one finite classical-output ensemble, exact upper bounds for all
three classes, and matching realizations in every dimension and along a
full-rank deformation.  Ohst et al.\ \cite[Theorem~24]{OZNPQ89} already
separate other strategy classes by channel-discrimination examples; the
present assertion does not claim a first memory advantage.  Its middle
optimization permits arbitrary unmeasured receiver storage and all
unit-reset feedback.  Lemma~\ref{lem:swapfilter90} is the step that prevents
a separable-tester upper from being incorrectly assigned to this larger
class.  These exact ensemble values and the local allocation orders
answer different questions; neither is substituted for the other.
